%% toc: P15 162 162.5,162.7,162.8,162.9
\chapter{Stopping at the First Variance Exceedance}
\label{sec:p15-w02c-p15-021-first-exceedance}

\Cref{thm:p15-finite-variation-convergence} localizes a square-integrable
martingale with the time \(T_c=\inf\{n\geq0:A_{n+1}>c\}\), which stops
one step before the predictable variance would pass the level \(c\).
That device is the one in the proof of Durrett's Theorem~4.5.2, where the
stopping time is \(N=\inf\{n:A_{n+1}>a^2\}\) and the bound used is
\(A_{N\wedge n}\leq a^2\) \cite[thm.~4.5.2, p.~240]{DurrettPTEAuthor2019}.
This chapter shows what is lost if one stops instead at the first time
the variance has already passed \(c\). The pathwise bound
\(A_{\cdot\wedge n}\leq c\) then holds only when stopping does nothing,
and a single martingale shows that the stopped processes need not be
bounded in \(L^2\) at any level, so the \(L^2\) convergence step of the
localization cannot be applied to them.

\emph{Inputs.} A real square-integrable martingale \(X\) with
\(X_0=0\), its predictable quadratic variation \(A\) as in
\cref{def:p15-predictable-variation}, and a level \(c>0\).
\emph{Conclusions.} The first-exceedance time is a stopping time equal
to \(T_c+1\); its stopped variation is bounded by \(c\) for every
horizon exactly on the event where the variance never passes \(c\);
with a deterministic bound on the variance increments it is bounded by
\(c\) plus that bound; without such a bound its second moments can be
unbounded at every level.

\section{The first-exceedance time}
\label{sec:p15-w02c-p15-021-definition}

\begin{hbdefinition}{def:p15-w02c-p15-021-first-exceedance}{The first-exceedance time of the variance process}
Let \(X\), \(a_n\) and \(A\) be as in
\cref{def:p15-predictable-variation}, and let \(c>0\). The
\emph{first-exceedance time} at level \(c\) is
\[
                       S_c=\inf\{n\geq0:A_n>c\},
\]
with \(\inf\emptyset=\infty\). The time \(T_c\) is the one of
\cref{thm:p15-finite-variation-convergence}.
\index{stopping time!first exceedance of the variance}
\end{hbdefinition}

\begin{hblemma}{thm:p15-w02c-p15-021-exceedance-shift}{proved=annex:p15-w02c-p15-021-exceedance-shift}{The first exceedance is one step after the localizing stop}
For every \(c>0\), \(S_c\geq1\), \(S_c\) is a stopping time,
\(S_c-1\) is \(T_c\), and \(\{S_c<\infty\}=\{A_\infty>c\}\).
On \(\{S_c<\infty\}\),
\[
              A_{S_c}=A_{T_c}+a_{S_c}>c,\qquad A_{T_c}\leq c.
\]
Here \(\infty-1=\infty\).
\end{hblemma}

\section{When the stopped bound holds}
\label{sec:p15-w02c-p15-021-bound}

\begin{hbtheorem}{thm:p15-w02c-p15-021-exceedance-bound}{proved=annex:p15-w02c-p15-021-exceedance-bound}{The first-exceedance bound holds only where stopping is idle}
For every \(c>0\) and every sample point,
\[
 A_{S_c\wedge n}\leq c\ \text{for every }n
 \quad\Longleftrightarrow\quad
 S_c=\infty
 \quad\Longleftrightarrow\quad
 A_\infty\leq c .
\]
Consequently the bound \(A_{S_c\wedge n}\leq c\) holds for every
\(n\), almost surely, exactly when \(A_\infty\leq c\) almost surely. In
that case \(X_{S_c\wedge n}=X_n\) for every \(n\), almost surely, so the
stop changes nothing.
If in addition \(a_n\leq K\) almost surely for every \(n\), with \(K\)
a deterministic constant, then
\[
          A_{S_c\wedge n}\leq c+K,\qquad
          EX_{S_c\wedge n}^2\leq c+K\qquad\text{for every }n.
\]
\end{hbtheorem}

\section{A martingale with unbounded first-exceedance stops}
\label{sec:p15-w02c-p15-021-example}

\begin{hbtheorem}{thm:p15-w02c-p15-021-unbounded-overshoot}{proved=annex:p15-w02c-p15-021-unbounded-overshoot}{A single variance jump defeats the first-exceedance localization}
Let \(J\) be a random variable with \(P(J=k)=2^{-k}\) for
\(k=1,2,\ldots\), and let \(\eta_1,\eta_2,\ldots\) be independent
signs with \(P(\eta_k=1)=P(\eta_k=-1)=1/2\), independent of \(J\). Put
\(\mathcal F_n=\sigma(J,\eta_1,\ldots,\eta_n)\) and
\[
 X_0=0,\qquad
 X_n=\sum_{k=1}^n2^{k/2}\mathbf1_{\{J=k\}}\eta_k\quad(n\geq1).
\]
Then \(X\) is a square-integrable martingale with \(X_0=0\),
\[
 a_n=2^n\mathbf1_{\{J=n\}},\qquad
 A_n=2^J\mathbf1_{\{J\leq n\}},\qquad
 A_\infty=2^J<\infty,\qquad
 EX_n^2=EA_n=n .
\]
For every \(c>0\):
\begin{enumerate}
\item \(X_{S_c\wedge n}=X_n\) and \(A_{S_c\wedge n}=A_n\) for every
\(n\), so \(\sup_nEX_{S_c\wedge n}^2=\infty\);
\item \(A_{S_c\wedge n}>c\) on the event \(\{J\leq n,\ 2^J>c\}\),
which has positive probability once \(2^n>c\), and
\(E\bigl((A_{S_c}-c)\mathbf1_{\{S_c<\infty\}}\bigr)=\infty\);
\item \(T_c=J-1\) on \(\{2^J>c\}\) and \(T_c=\infty\) on
\(\{2^J\leq c\}\), \(X_{T_c\wedge n}=X_n\mathbf1_{\{2^J\leq c\}}\) and
\(EX_{T_c\wedge n}^2=EA_{T_c\wedge n}\leq c\) for every \(n\).
\end{enumerate}
\end{hbtheorem}

\begin{hbwarning}{def:p15-w02c-p15-021-what-fails}{What the example does and does not show}
In \cref{thm:p15-w02c-p15-021-unbounded-overshoot} the path \(X_n\) is
constant from time \(J\) on, so it converges, and
\cref{thm:p15-finite-variation-convergence} is not contradicted. What
fails is the route: stopped at the first exceedance, the martingale is
not bounded in \(L^2\) at any level \(c\), so neither
\cref{thm:p15-quadratic-maximal-bound} nor the \(L^2\) convergence
theorem applies to the stopped process. The reason is the overshoot
\(a_{S_c}=A_{S_c}-A_{T_c}\), the one variance increment that the
first-exceedance time includes and the localizing time \(T_c\) excludes.
Because \(a_{S_c}\) is known at time \(T_c=S_c-1\), stopping at \(T_c\)
uses no information that is not already available.
\end{hbwarning}

\section{Relationships and sources}
\label{sec:p15-w02c-p15-021-relations}

The same one-step-early stop appears in the square-root estimate
\cref{thm:p15-square-root-variation-bound}, whose proof stops at
\(T_{a^2}\). For a martingale, rather than its variance, exits by a
level are controlled through a deterministic bound on the increments in
\cref{thm:p15-bounded-jump-variation-dichotomy}; the bounded-increment
part of \cref{thm:p15-w02c-p15-021-exceedance-bound} is the analogous
statement for the variance process. Durrett states the convergence on
\(\{A_\infty<\infty\}\) as Theorem~4.5.2 and proves it with the time
\(N=\inf\{n:A_{n+1}>a^2\}\), noting that \(A_{n+1}\) is
\(\mathcal F_n\)-measurable \cite[thm.~4.5.2, p.~240]{DurrettPTEAuthor2019}.
The characterization of the first-exceedance bound, its bounded-increment
form and the example are proved here; they are not stated in that source.
